% !TeX root = ../papers/nc-p3.tex
\section{階数15の全通常点を分離する三つの添字保持不変量}
\label{sec:rank-fifteen-global-invariants}
\begin{theorem}\label{thm:rank-fifteen-global-invariants}
正の固定曲線の任意の二つの通常な階数15窓が添字保持同型である必要十分条件は、
前節の $(x,p,q)$ が一致することである。二積の生成次元に条件は要らない。
特にこの一判定は分岐、座標端の例外水準、および異なる倍率間の比較を含む。
\end{theorem}
\begin{proof}
十分性は定理\ref{thm:rank-fifteen-rational-quotient}による。必要性のため
$H\in M_4\otimes M_4$ を任意の基底で規格化した反対切断とする。
$T$ をこれに対応する線型写像、$S_1=T^*T$, $S_0=TT^*$ とし、
$\Gamma$ を第一行列因子の転置とする。$D(H)$ は
節\ref{sec:rank-fifteen-endpoint-orbits}の二置換による五重縮約、
$B(S_j)$ は命題\ref{prop:rank-fifteen-oriented-cubic}の対称三次形式の縮約分子である。
次の量は全て元のテンソルから計算できる。
\[
 A_2=\frac{\operatorname{tr}H^2}{16},\qquad
 A_3=\frac{\operatorname{tr}H^3}{96},\qquad
 A_4=\frac{A_2^2-\operatorname{tr}(S_1^2)/64}{3}.
\]
すると前節の座標は次の式を持つ。
\[
\begin{aligned}
 x&=A_4 A_2/A_3^2, &y&=A_4^3/A_3^4,\\
 p&=-\frac{A_4\bigl(D(H)-D(\Gamma H)\bigr)}{128 A_3^3},
 &q&=\frac{B(S_1)-B(S_0)}{3072 A_3^2}.
\end{aligned}
\]
独立共役で全ての縮約は不変である。転置を取った因子の共役は逆転置行列の共役に移るので
$D(\Gamma H)$ も不変である。各分数の分子と分母の $H$ に関する次数は等しい。
規格化した係数表示では $A_3=U-1$ であり、階数15の全通常点で単元である。

恒等式を検算する。前節の斉次八量で $A_2=P_2$, $A_3=P_3$, $A_4=h^2$ となる。
実際、$P_j$ をここではラベル $j$ のPauli基底テンソルとすると、
$P_j^2=(-1)^{e(j,j)}I$ である。
$\mu_j=(-1)^{e(j,j)}\gamma_j^2$ と置けば、$S_1,S_0$ は対応する各局所Pauli基底で固有値
$16\mu_j$ を持つ。これは $T$ が一つのPauli行列を他方の行列へ、そのトレース双対による
係数を掛けて送ることから従う。よって $\operatorname{tr}H^2/16=\sum_j\mu_j$、
$\operatorname{tr}S_1^2/256=\sum_j\mu_j^2$ である。$H^3$ のトレースには
$j,k,j+k$ の組だけが寄与し、両因子の積の符号を掛ける。

$B(S_1)$ では入力側で可換な非零三ラベル $j,k,j+k$ の順序対 $(j,k)$ ごとに、
三つの固有値の積が一回寄与する。反可換な対の寄与は零である。
実際、局所Pauli基底 $E_j$ のトレース双対は $E_j/(4E_j^2)$ であり、可換な場合は
$d(E_j,E_k,E_{j+k})$ とその双対での値の積が1になる。
$B(S_0)$ では出力側の可換性を使う。この差の和を取ると $3072E$ となる。
五重縮約では入力の二閉路によりラベル列が $(a,c,b,a+b,c)$ となる。
第一因子の転置で符号が変わる項に限り、両因子で各閉路の積の符号を掛ける。
三閉路のトレース係数 $4^3$ と、差を取る因子2を戻すと
$D(H)-D(\Gamma H)=-128hD_3$ を得る。
\path{computations/verify_rank_fifteen_global_invariants.py} はこれらの多項式恒等式を
五つの二次式で簡約して検算する。さらに101上の四つの枠ではGram表示から
$H,T$ を再構成し、閉路の直接縮約とも比較する。
これで四つの分数が前節の値と一致し、全点での必要性が従う。
\end{proof}

\begin{corollary}
通常な階数15窓が語反転で得る反対窓と添字保持同型となるのは分岐上に限り、
一位置移動した反対窓と添字保持同型となるのは $U=9$ のときに限る。
二つの軌跡に共通の通常点はない。
\end{corollary}
\begin{proof}
前節の二操作で反転される座標はそれぞれ $q,p$ であり、
規格化テンソルでの群軌道の一致が十分である。上の縮約は、変換先の倍率が非通常でも
そのテンソルに直接適用できる。通常域では $q=0$ は $2n^2+1-U=0$ すなわち分岐で、
$p=0$ は $U=9$ である（$zkm\ne0$）。
両方が成り立てば $K=\pm3$, $H=\pm4$ なので $K+H,K-H$ の一方が $\pm1$ となる。
すると $z=1$ または $z\lambda^2=1$ となり通常条件に反する。
\end{proof}
